\documentclass[10pt,a4paper]{amsart}
\usepackage{amsmath,amssymb,amsthm,graphicx}
\usepackage{hyperref}
\usepackage{tikz}
%\usepackage{xcolor}
\usepackage{pict2e}
\let\arc\relax
%\usepackage{harpoon}
\usepackage{curves}
%\usepackage{chngcntr}
%\usepackage{apptools}
\usepackage{enumerate}

%\usepackage{verbatim}
%\usepackage[nospace,noadjust]{cite}
\usepackage{accents}
\usetikzlibrary{calc,decorations.markings}
\unitlength1cm

%%%%%%%%%%%%%%FONTS%%%%%%%%%%%%%%%%%%%%%%%%%%%%%%
\newcommand{\R}{\mathbb{R}}
\newcommand{\Z}{\mathbb{Z}}
\newcommand{\N}{\mathbb{N}}
\newcommand{\C}{\mathbb{C}}
\newcommand{\Id}{I}

%%%%%%%%%%%%%%MATHCAL%%%%%%%%%%%%%%%%%%%%%%%
\newcommand{\cK}{\mathcal{K}}
\newcommand{\cE}{\mathcal{E}}
\newcommand{\cR}{\mathcal{R}}
\newcommand{\cW}{\mathcal{W}}

%%%%%%%%%%%%%%%%%%ABBRS%%%%%%%%%%%%%%%%%%%%%%%%%%%%%

%%%%%%% OBJECTS SPECIFIC TO THIS PAPER%%%%%%%%%

\newcommand{\hw}{\widehat w}
\newcommand{\tw}{\widetilde w}

\newcommand{\bQ}{ Q}
\newcommand{\hQ}{{\widehat Q}}
\newcommand{\tQ}{{\widetilde Q}}

\newcommand{\sbQ}{Q^\star}
\newcommand{\shQ}{{\widehat Q^\star}}
\newcommand{\stQ}{{\widetilde Q^\star}}

\newcommand{\smu}{\mu^\star}
\newcommand{\shmu}{\widehat \mu^\star}
\newcommand{\stmu}{\widetilde \mu^\star}

\newcommand{\bmu}{{\mu}}
\newcommand{\hmu}{{\widehat \mu}}
\newcommand{\tmu}{{\widetilde \mu}}


\newcommand{\bP}{P}
\newcommand{\hP}{{\widehat P}}
\newcommand{\tP}{{\widetilde P}}

\newcommand{\sbP}{P^\star}
\newcommand{\shP}{{\widehat P^\star}}
\newcommand{\stP}{{\widetilde P^\star}}

\newcommand{\hR}{{\widehat R}}

\newcommand{\bv}{v}
\newcommand{\hv}{{\widehat v}}
\newcommand{\tv}{{\widetilde v}}

\newcommand{\sbv}{v^\star}
\newcommand{\shv}{{\widehat v^\star}}
\newcommand{\stv}{{\widetilde v^\star}}

\newcommand{\stu}{{\widetilde u^\star}}

\newcommand{\hF}{\widehat F}
\newcommand{\hW}{\widehat W}

\newcommand{\Fw}{\frac{F^2}{w}}
\newcommand{\hFw}{\frac{\hF^2}{\hw}}

\newcommand{\Cau}{\mathcal C_{\Sigma}}
\newcommand{\Cauv}{\mathcal C_{v}}
\newcommand{\Cauhv}{\mathcal C_{\hv}}
\newcommand{\Cautv}{\mathcal C_{\tv}}

\newcommand{\sCau}{\mathcal C_{\Sigma^\star}}
\newcommand{\sCauv}{\mathcal C_{\bv^\star}}
\newcommand{\sCauhv}{\mathcal C_{\hv^\star}}
\newcommand{\sCautv}{\mathcal C_{\stv}}

\newcommand{\Ltwo}{L^2(\Sigma)}
\newcommand{\sLtwo}{L^2(\Sigma^\star)}

\newcommand{\Sigl}{\Sigma^\ell}
\newcommand{\Sigr}{\Sigma^r}

%%%%%%%%%%%%%%%%OTHER SHORTCUTS%%%%%%%%%%%
\newcommand{\nn}{\nonumber}
\newcommand{\ul}{\underline}
\newcommand{\ol}{\overline}
\newcommand{\pa}{\partial}
\newcommand{\ti}{\tilde}
\newcommand{\wti}{\widetilde}
\newcommand{\wha}{\widehat}
\newcommand{\norm}[1]{\lVert#1 \rVert}
\newcommand{\abs}[1]{\lvert#1 \rvert}
%\newcommand{\Id}{\mathbb{I}}
\newcommand{\I}{\mathrm{i}}
\newcommand{\E}{\mathrm{e}}
\newcommand{\re}{\mathop{\mathrm{Re}}}
\newcommand{\im}{\mathop{\mathrm{Im}}}
\newcommand{\eps}{\varepsilon}
\newcommand{\dist}{\mathop{\mathrm{dist}}}

\newtheorem{Theorem}{Theorem}[section]
\newtheorem{Corollary}[Theorem]{Corollary}
\newtheorem{Lemma}[Theorem]{Lemma}
\newtheorem{Proposition}[Theorem]{Proposition}
 { \theoremstyle{definition}
\newtheorem{Definition}[Theorem]{Definition}
\newtheorem{Note}[Theorem]{Note}
\newtheorem{Example}[Theorem]{Example}
\newtheorem{Remark}[Theorem]{Remark} }

\numberwithin{equation}{section}
\usepackage{cleveref}
\crefname{originaltheorem}{original theorem}{original theorems}
\Crefname{originaltheorem}{Original theorem}{Original theorems}
\makeatletter
\newlabel{TheoremInh}{{internal}{source}{}{sourcepointer-1}{}}
\newlabel{TheoremInh@cref}{{[originaltheorem][0][]internal}{[1][0][]source}{}{}{}}
\makeatother
\crefname{originalcorollary}{original corollary}{original corollarys}
\Crefname{originalcorollary}{Original corollary}{Original corollarys}
\makeatletter
\newlabel{CorHn}{{internal}{source}{}{sourcepointer-2}{}}
\newlabel{CorHn@cref}{{[originalcorollary][0][]internal}{[1][0][]source}{}{}{}}
\makeatother
\crefname{originalequation}{original equation}{original equations}
\Crefname{originalequation}{Original equation}{Original equations}
\makeatletter
\newlabel{LegRHP}{{internal}{source}{}{sourcepointer-3}{}}
\newlabel{LegRHP@cref}{{[originalequation][0][]internal}{[1][0][]source}{}{}{}}
\makeatother
\crefname{originaltheorem}{original theorem}{original theorems}
\Crefname{originaltheorem}{Original theorem}{Original theorems}
\makeatletter
\newlabel{UnifLeg}{{internal}{source}{}{sourcepointer-4}{}}
\newlabel{UnifLeg@cref}{{[originaltheorem][0][]internal}{[1][0][]source}{}{}{}}
\makeatother
\title{The explicit form of the inverse of the Cauchy type operator}
\author{Deift, Percy, Piorkowski, Mateusz}
\date{doi:10.3842/SIGMA.2024.004; CC BY 4.0}
\begin{document}
\maketitle
\setcounter{section}{4}
\setcounter{Theorem}{5}
\noindent\textit{Editorial note.} All delivered mathematical text is the author's own, copied from the pinned SIGMA article source; no statement or proof step was rewritten, reordered or removed. The counted statement is the labeled proposition environment at source lines 1139-1145 and the counted proof is the proof environment at source lines 1146-1186.\par\medskip
\noindent\textit{Editorial note.} The intake census pairs this statement with a proof environment that the source places after two further statements of the same block; reading the source shows that the proof derives exactly the formula of the counted statement and that the two statements placed before it are quoted results of the cited literature, so the counted statement has its own proof.\par\medskip
\noindent\textit{Editorial note.} The proof cites published references of the author's bibliography by key; the cited entries are copied line for line from the author's own bibliography into a retained bibliography fragment, and hyperref is added because the publisher class, not the author preamble, supplied the link command.\par\medskip
\noindent\textit{Editorial note.} The retained preamble of this unit was generated mechanically for this delivery: it copies the author's own package, numbering, macro and theorem declaration block from the pinned source (source lines 9-133), with the publisher class and page setup lines removed.\par\medskip
\noindent\textit{Allowed standard prerequisites (editorial).}\begin{itemize}
\item The equivalence between invertibility of the singular operator and unique solvability of the inhomogeneous Riemann-Hilbert problem. same article, the theorem at source lines 1128-1130. Converts the equation for the unknown into an inhomogeneous Riemann-Hilbert problem whose solution can be written down.
\item The Legendre Riemann-Hilbert problem and the uniqueness of its solution. same article, the problem at source line 996 and the statement at source line 1115. Fixes the normalization and the jump of the solution used to express the inverse.
\end{itemize}\medskip
\ifdefined\hypertarget\hypertarget{sourcepointer-1}{}\fi
\noindent\textit{Source reference.} Original theorem internal: sources/sigma-2024-004-recurrence-coefficients-for-orthogonal-polyn/pinned-original-source.tex, source lines 1128--1130. The counted proof refers to that the author's theorem relating invertibility to unique solvability of the inhomogeneous problem; it is the author's own earlier display, definition, figure or statement in the same article and the whole argument of the counted span is retained verbatim.\par\medskip
\ifdefined\hypertarget\hypertarget{sourcepointer-2}{}\fi
\noindent\textit{Source reference.} Original corollary internal: sources/sigma-2024-004-recurrence-coefficients-for-orthogonal-polyn/pinned-original-source.tex, source lines 1135--1137. The counted proof refers to that the author's corollary on the uniform invertibility of the operators along a sequence; it is the author's own earlier display, definition, figure or statement in the same article and the whole argument of the counted span is retained verbatim.\par\medskip
\ifdefined\hypertarget\hypertarget{sourcepointer-3}{}\fi
\noindent\textit{Source reference.} Original equation internal: sources/sigma-2024-004-recurrence-coefficients-for-orthogonal-polyn/pinned-original-source.tex, source lines 996--1000. The counted proof refers to that the author's displayed Legendre Riemann-Hilbert problem; it is the author's own earlier display, definition, figure or statement in the same article and the whole argument of the counted span is retained verbatim.\par\medskip
\ifdefined\hypertarget\hypertarget{sourcepointer-4}{}\fi
\noindent\textit{Source reference.} Original theorem internal: sources/sigma-2024-004-recurrence-coefficients-for-orthogonal-polyn/pinned-original-source.tex, source lines 1115--1119. The counted proof refers to that the author's statement on the uniform behaviour of the Legendre problem; it is the author's own earlier display, definition, figure or statement in the same article and the whole argument of the counted span is retained verbatim.\par\medskip
\begin{Proposition}\label{ExpForm}
The inverse of $1-\Cautv$ has the explicit form
\begin{align}\label{ExpInv}
 (1-\Cautv)^{-1}\colon\ \Ltwo \to \Ltwo, \qquad g \mapsto g + \Cau^-\big(g (\tv-\Id) \tQ_+^{-1}\big)\tQ_-.
\end{align}

\end{Proposition}

\begin{proof}
From Theorem \ref{TheoremInh}, we know that for $g \in \Ltwo$, the (unique) solvability of the equation
\begin{align}\label{SingIE}
(1-\Cautv)\psi = g, \quad g\in \Ltwo,
\end{align}
in $\Ltwo$ is equivalent to the (unique) solvability of the following inhomogeneous RH problem.

\subsection*{Inhomogeneous Legendre RH problem}
For a given $g \in \Ltwo$, one seeks an $f \in \Ltwo$, such that $m_\pm = \Cau^\pm(f)+g$ satisfies the jump relation:
\begin{align*}
 m_+(s) &= m_-(s) \tv(s), \qquad s \in \Sigma.
\end{align*}

Note that Theorem \ref{UnifLeg} together with Corollary~\ref{CorHn} imply that there exists a constant $c$ independent of $n\in \mathbb{N}$ and $g\in \Ltwo$ such that the inhomogeneous Legendre RH problem has a unique solution $m_\pm$ with $\Vert m_- \Vert_{\Ltwo} \leq c\Vert g \Vert_{\Ltwo}$. We shall briefly recall the exact relation between \eqref{SingIE} and the inhomogeneous Legendre RH problem (cf.\ \cite[Section~2]{DZSob}), as it will be needed later in the proof. First, note that if we have a solution $m_\pm$ to the inhomogeneous Legendre RH problem, we must have $f = m_+ - m_- = m_-(\tv-\Id)$. If we now set $\psi := m_- = \Cau^-(f)+g \in \Ltwo$, it follows that
\begin{align*}
 (1-\Cautv)\psi &= \psi - \Cau^-(\psi (\tv-\Id))= \Cau^-(f) + g - \Cau^-(m_-(\tv-\Id)))= g.
\end{align*}
On the other hand, having a solution $\psi$ to the integral equation \eqref{SingIE}, we can define $m_\pm := \Cau^\pm(\psi (\tv-\Id)) + g$ and compute, using the Sokhotski--Plemelj (SP) formula,
\begin{align}\label{proofm}
 m_+ &= \Cau^+(\psi (\tv-\Id)) + g\overset{\text{SP}}{=} \psi (\tv-\Id) + \Cautv(\psi) + g= \psi \tv= m_- \tv,
\end{align}
as $m_- = \Cautv \psi + g = \psi$.
We can now derive an expression for the operator $(1 - \Cautv)^{-1}$. Assume $g \in \Ltwo$ is given and let $\psi \in \Ltwo$ be the unique solution of $(1-\Cautv)\psi = g$. Then $m_\pm = \Cau^\pm(f)+g$ with $f = \psi (\tv-\Id)$ solves the corresponding inhomogeneous RH problem, as we have seen in equation~\eqref{proofm}. We want to find an expression for $\psi = m_-$ in terms of $g$. To derive \eqref{ExpInv}, we start with
\begin{align}
& m_+ = m_- \tv, \qquad
 \Cau^+(f) + g = (\Cau^-(f) + g)\tilde v,\qquad \big(\Cau^+(f) + g\big)\tQ_+^{-1} = (\Cau^-(f) + g)\tQ_-^{-1},\nonumber
 \\
& \Cau^+(f)\tQ_+^{-1} - \Cau^-(f)\tQ_-^{-1} = g\big(\tQ_-^{-1}-\tQ_+^{-1}\big).\label{derInv}
 \end{align}
Observe that the left- and right-hand sides in the last line might not lie in $L^2(\Sigma)$. However, using property (iv) of $\tQ$ from the RH problem (see Section~\ref{LegRHP}), we see that they lie in $L^{2-\epsilon}(\Sigma)$ for any $\epsilon \in (0,1)$. Hence, if we define $H = \Cau(f) \tQ^{-1}$, we see that $H$ is analytic in $\C \setminus \Sigma$, satisfies $H_\pm \in L^{2-\epsilon}(\Sigma)$ and vanishes at $\infty$. By the Sokhotski--Plemelj formula, it follows that $H = \Cau(H_+-H_-)$. Thus, applying $\Cau$, which is understood to act on the space $L^{2-\epsilon}(\Sigma)$, in the last line of \eqref{derInv}, we obtain
 \begin{align*}
& \Cau(f) \tQ^{-1} = \Cau\big(\underbrace{g\big(\tQ_-^{-1}-\tQ_+^{-1}\big)}_{g(\tilde v-\Id) \tQ_+^{-1}}\big),\qquad
 \Cau(f) = \Cau\big(g (\tv-\Id) \tQ_+^{-1}\big)\tQ,
 \\
& \Cau^-(f) + g = g + \Cau^{-}\big(g (\tv-\Id) \tQ_+^{-1}\big)\tQ_-,\qquad
 \psi = g + \Cau^-\big(g (\tv-\Id) \tQ_+^{-1}\big)\tQ_-.
\end{align*}
Note here that $\Cau^-\big(g (\tv-\Id) \tQ_+^{-1}\big)\tQ_-$ is a priori a function in $L^{2-\epsilon}(\Sigma)$, however as $\psi = (1-\Cautv)^{-1}g \in L^2(\Sigma)$ by Theorem \ref{UnifLeg}, we conclude that indeed $\Cau^-\big(g (\tv-\Id) \tQ_+^{-1}\big)\tQ_- = \psi - g \in \Ltwo$.

We have thus proved that $(1-\Cautv)^{-1}$ must indeed have the formed stated in \eqref{ExpInv}.
\end{proof}

\begin{thebibliography}{99}
\bibitem{DZSob}
Deift P., Zhou X., Long-time asymptotics for solutions of the {NLS} equation
 with initial data in a weighted {S}obolev space, \href{https://doi.org/10.1002/cpa.3034}{\textit{Comm. Pure Appl.
 Math.}} \textbf{56} (2003), 1029--1077, \href{https://arxiv.org/abs/math.AP/0206222}{arXiv:math.AP/0206222}.
\end{thebibliography}
\end{document}
